% RH paper — § Introduction
% Drafted body for the introduction. This section is framing only and
% forward-references the definitions and theorem statements resolved in
% §§3--7.
%
% statements-of-record rows resolved here: none (this section is
% framing only; theorem/definition rows resolve in §3–§7 per
% paper/notes/statements-of-record.yml target_section_hint fields).

\section{Introduction}
\label{sec:intro}

What would it take to settle whether all nontrivial zeros of the
Riemann zeta function lie on the critical line?  In its standard
form, the Riemann hypothesis asserts that every nontrivial zero
\(\rho\) of \(\zeta\) satisfies
\[
  \Real(\rho)=\frac12.
\]
Classical attacks on this statement seek an unconditional proof inside
the analytic-number-theoretic substrate itself: spectral Hilbert--Pólya
programs, adelic and noncommutative-geometric approaches, Hilbert
spaces of entire functions, Weil-positivity formulations, and related
programs all work within their own classical carriers.  This paper
takes a different route.  It asks whether RH can be read as the
vanishing of a positive typed quantity and then made to vanish by a
structural closure mechanism.  The quantity is the anti-invariant zero
ledger
\[
  \AZ=\sum_{\rho\in \Znt} m_\rho\,
       \left|\Real(\rho)-\frac12\right|^2,
\]
where \(\Znt\) is the multiplicity-weighted ledger of nontrivial zeros
of the completed zeta function.  Its vanishing is, by a direct
positive-sum calculation, equivalent to the critical-line statement.
The remaining question is not how to prove RH by one of the classical
routes above, but how a formed closure can force \(\AZ\) to vanish.

The framework used here is Six Birds Theory, a typed and audited
calculus for naming closure formation in mathematical structures.  A
formed closure, in the sense used in this paper, is a typed presentation
of a mathematical substrate together with the assumption that its
carriers, observables, comparison maps, and audit data have been
admissibly closed under the framework's formation discipline.  The
substrate for the present paper is the saturated completed Selberg trace
closure \(\Sel\): a typed shell aggregating the completed zeta data, a
lawful trace instrument \(\Itr\), the functional-equation symmetry, the
nontrivial-zero ledger, the anti-invariant zero ledger, visibility data,
and audit provenance.  The functional-equation involution is
\[
  \JL(s)=1-\bar{s},
\]
whose fixed locus is the critical line.  The corresponding
anti-invariant readout is the signed real-part displacement
\[
  \psimRH(s)=\Real(s)-\frac12.
\]
These objects are defined in \Cref{sec:involution_and_ledger}; the
shell \(\Sel\) is defined in \Cref{sec:sat_sel_shell}.

The headline theorem is conditional.  Under the standard Six Birds
closure assumption on \(\Sel\), and given the recognition source
\(\GamSDTCSelberg\), the Riemann hypothesis holds on the typed zero
ledger of \(\Sel\).  Here a recognition source means a named structural
hypothesis supplied as content of the formed layer, not a theorem
derived from lower framework primitives.  In this instance,
\(\GamSDTCSelberg\) is the Selberg-class instance of Self-Dual Trace
Confinement from the sibling duality-confinement paper
\cite{TsiokosSDTC2026}.  It supplies, on the formed Selberg trace
layer, a sequence of completed domination records
\[
  \AZ \preceq B_n,
  \qquad
  \trace B_n \to 0.
\]
The duality-confinement membrane theorem of the sibling paper then
forces \(\AZ=0\) in the typed positive cone, and the translation
theorem of the present paper converts that vanishing into
\(\Real(\rho)=1/2\) for every \(\rho\in\Znt\).  Outside the Six Birds
framing, the result should therefore be read exactly as the conditional
theorem
\[
  \GamSDTCSelberg \Longrightarrow \mathrm{RH},
\]
not as an unconditional standard-ZFC proof of RH.

The translation step is deliberately elementary.  \Cref{sec:translation_theorem}
proves Theorem T:
\[
  \AZ=0
  \quad\Longleftrightarrow\quad
  \forall \rho\in\Znt,\ \Real(\rho)=\frac12.
\]
The forward direction uses only that each summand
\(m_\rho |\Real(\rho)-1/2|^2\) is nonnegative and that
\(m_\rho>0\).  If the total sum is zero, every squared distance from
the critical line is zero.  The reverse direction is substitution:
if every typed zero has real part \(1/2\), then every summand
vanishes, so \(\AZ=0\).  This theorem is proven by direct positivity
argument: it does not carry the recognition-source burden.
Rather, it cleanly translates the positive ledger statement produced
by the duality-confinement mechanism into the number-theoretic
critical-line statement.

The paper's scope is correspondingly narrow and explicit.  First, the
main result is conditional on \(\GamSDTCSelberg\).  We do not claim
unconditional RH in standard ZFC, and we do not claim that
\(\GamSDTCSelberg\) is derivable from framework primitives.  The
development history includes a three-route derivation sweep whose role
is due diligence: it supports the recognition-source status by
documenting that the domination content is not obtained from the lower
primitive vocabulary alone.  Second, the recognition source is not
installed as a Lean axiom.  It is represented in the formalization as a
typed structure carrier whose fields record the required domination
data; the project-wide source rule bans declarations using the tokens
\texttt{axiom}, \texttt{opaque}, \texttt{constant}, \texttt{sorry}, and
\texttt{admit}.  The conditional theorem takes a value of that carrier
as an explicit hypothesis.

Third, the zero ledger is typed.  The formal encoding parameterizes the
real-part coordinate by a typed real-coordinate carrier, and the
identification of that coordinate with the usual real part of a
classical nontrivial zero of \(\zeta\) is made by construction-parameter
assignment when the shell is instantiated.  This is the correct
mathlib-free abstraction for the present formalization; a Lean-side
derivation of the analytic identification with the classical zero set
is left as a future extension.  Fourth, admissibility of \(\Sel\) under
the seven Foundations-II schemas is carried as an audit/admissibility
witness in the shell, not derived in this paper.  Finally, the
classical RH programs named above are neither refuted nor circumvented
in their native sense.  They aim at unconditional RH inside classical
analytic number theory; the present paper establishes a conditional
formed-layer result on the typed Selberg trace shell.  The generalized
Riemann hypothesis is also outside the scope of this paper: we
instantiate the construction for \(L=\zeta\) only.

The remainder of the paper follows this division of labor.
\Cref{sec:framework} gives the minimal Six Birds background needed for
the closure reading.  \Cref{sec:involution_and_ledger} constructs the
functional-equation involution, the anti-invariant readout, the
nontrivial-zero ledger, and \(\AZ\).  \Cref{sec:sat_sel_shell}
packages the saturated Selberg trace closure.  \Cref{sec:translation_theorem}
proves Theorem T.  \Cref{sec:recognition_source} states the recognition
source and its provenance, and \Cref{sec:landing_chain} composes it
with the duality-confinement master theorem and Theorem T to obtain the
conditional landing chain.  \Cref{sec:scope_and_nonclaims} records the
nonclaims explicitly, and \Cref{sec:discussion,sec:conclusion} place
the result in the broader structural program.  The formalization
appendix records the ten mechanized RH entries with faithful semantic
alignment, together with the one recognition-source carrier and the
representation notes needed to read the Lean coverage correctly.
